
%%%%%%%%%%  scan idx 36  |  folio 25  %%%%%%%%%%


\underline{Exposición II}


\underline{\underline{Propiedades de finitud del grupo fundamental}}


\underline{por Michèle RAYNAUD}

Sean  $k$  un cuerpo separablemente cerrado de característica  $p \geqslant 0$ , $X$  un esquema conexo de tipo finito sobre  $k$  y  $\Pi_1^{(p)}(X)$  el máximo cociente «primo a $p$» de  $\Pi_1(X)$ . En esta exposición se demuestra que, si  $X$  es una superficie, $\Pi_1^{(p)}(X)$  es topológicamente de presentación finita, y que lo mismo ocurre para un  $X$  arbitrario si se admite la resolución de singularidades.

En el caso de una superficie, el método de demostración, debido a J.-P. Murre, consiste en reducirse al caso en que se tiene una fibración sobre  $P^1$  con las propiedades (i), (ii) y (iii) de la proposición siguiente.

\underline{Proposición} 2.1.  \underline{Sean}  $k$  \underline{un cuerpo algebraicamente cerrado} y  $S$  \underline{una superficie proyectiva}, \underline{irreducible} y \underline{lisa sobre}  $k$ . \underline{Entonces se puede encontrar una explosión}

\[ g : S' \longrightarrow S \quad , \]
\underline{donde}  $S'$  \underline{es una superficie regular}, \underline{provista de una fibración}

\[ f : S' \longrightarrow P^1 \quad , \]
\underline{de modo que}  $f$  \underline{tenga una sección}  $\sigma$  \underline{y se satisfagan las condiciones siguientes}:

(i)   $f$  \underline{es liso en los puntos de}  $\sigma(P^1)$ .

(ii)  \underline{las fibras de}  $f$  \underline{son curvas geométricamente íntegras que solo tienen singularidades ordinarias.}

(iii) \underline{existe un abierto no vacío}  $U$  \underline{de}  $P^1$  \underline{tal que las fibras de}  $f$  \underline{en los puntos de}  $U$  \underline{son lisas.}

La proposición se deduce de la existencia de una inmersión proyectiva  $S \longrightarrow P^r$  respecto de la cual se puede encontrar un haz de Lefschetz (XVII 2.5). Sea, en efecto,  $D = \{H_t\}_{t \in P^1}$  tal haz de hiperplanos; recordemos que se
%%%%%%%%%%  scan idx 37  |  folio 26  %%%%%%%%%%
tienen las propiedades siguientes (loc. cit. 2.2):

a) el eje  $\Delta$  de  $D$  corta transversalmente a  $S$ ; el subesquema cerrado  $\Delta \cap S$  es, por tanto, reducido y está formado por un conjunto finito de puntos cerrados de  $S$ , $x_1,...,x_n$ .

b) existe un abierto no vacío  $U$  de  $P^1$  tal que, para todo punto  $t \in U$ , $H_t$  corta transversalmente a  $S$ .

c) para todo punto  $t \in P^1-U$ , $H_t$  corta transversalmente a  $S$  salvo en un punto, que es un punto singular cuadrático ordinario.

Por cada punto  $x$  de  $S$  que no pertenece a  $\Delta$  pasa un único hiperplano  $H_t$ , y la aplicación que asocia a  $x$  el punto  $t \in P^1$  es una aplicación racional  $a$  de  $S$  en  $P^1$  cuyo dominio de definición es  $S-(\Delta \cap S)$ .

Consideremos la explosión  $g : S' \longrightarrow S$  a lo largo del subesquema cerrado  $S \cap \Delta$  de  $S$ . El esquema  $S'$  es regular y la fibra  $S'_{x_i}$  sobre un punto  $x_i$  se identifica con el fibrado proyectivo definido por el espacio tangente a  $S$  en  $x_i$ , por lo que es isomorfa a  $P^1$  (EGA IV 19.4.2, 19.4.4). Además, la aplicación racional  $f = ag$  está definida en todas partes (XVIII 3.1). Finalmente, a cada punto  $x_i$  se le puede asociar la sección de  $f$  que hace corresponder a un punto  $t \in P^1$  el vector tangente en  $x_i$  a  $H_t \cap S$ .

Comprobemos las condiciones (i), (ii) y (iii). Obsérvese que, para cada  $t \in P^1$ , las curvas  $S'_t$  y  $S_t = H_t \cap S$  son isomorfas; en efecto, el morfismo  $S'_t \longrightarrow S_t$  es un isomorfismo salvo quizá en los puntos  $x_i$ , puntos en los que  $S_t$  es regular; además, como  $S'_t$  es una intersección completa en  $S'$ , $S'_t$  es íntegra y, por tanto, isomorfa a  $S_t$ .

\sisi{Las condiciones (i) y (ii) no son sino las condiciones a) y b).}{Las propiedades de las fibras enunciadas en (ii) y (iii) resultan de las condiciones a), b) y c).} Finalmente,  $f$  es liso en los puntos de  $\sigma$ , pues  $f$  es plano y las fibras  $S'_t$  son geométricamente regulares en los puntos  $\sigma(t)$ .


%%%%%%%%%%  scan idx 38  |  folio 27  %%%%%%%%%%

\underline{Corolario} 2.2. \underline{Con las notaciones de} 2.1, \underline{sea} $Y$ \underline{un divisor de cruces normales en} $S$ , $Y' = Y \times_S S'$ . \underline{Entonces se pueden elegir} $f$ , $g$ , $U$ \underline{de modo que} $Y'\times_{P^1}U$ \underline{sea un divisor de cruces normales relativo a} $U$ .

Se sabe que, si existe un haz de Lefschetz para una inmersión $S \longrightarrow P^r$ , los haces de Lefschetz forman un abierto de la variedad de rectas de $P^r$ (XVII 3.2.1). Lo mismo ocurre con el conjunto de los haces de Lefschetz cuyo eje no corta a $Y$ y para los cuales existe un abierto no vacío $U_1$ de $P^1$ tal que la intersección de $Y$ y $H_t$ sea lisa para $t \in U_1$ . Basta entonces construir $f$ y $g$ a partir de un haz que satisfaga estas dos condiciones y sustituir el abierto considerado por $U \cap U_1$ .

2.3.0. Sea $G$ un grupo profinito. Recordemos que se dice que $G$ es topológicamente de generación finita si existe un entero $n$ tal que $G$ es isomorfo a un cociente del grupo profinito libre con $n$ generadores $F(n)$ ; sea $K = Ker(F(n) \longrightarrow G)$ ; se dice que $G$ es topológicamente de presentación finita si, además, se puede encontrar un número finito de elementos $k_1,...,k_r$ de $K$ tales que el menor subgrupo normal cerrado de $K$ que contiene a $k_1,...,k_r$ es igual a $K$ .

\underline{Teorema} 2.3.1. \underline{Sean} $k$ \underline{un cuerpo separablemente cerrado de característica} $p \geqslant 0$  y $X$ \underline{un esquema conexo de tipo finito sobre} $k$ . \underline{Supóngase que los esquemas de tipo finito y de dimensión} $\leqslant dim(X)$ \underline{sobre una clausura algebraica de} $k$ \underline{admiten una desingularización fuerte} (SGA 5 I 3.1.5) (\underline{condición que se cumple si} $dim X = 2$ \underline{según} [1] \underline{o si} $\sisi{k}{\operatorname{char}k} = 0$ \underline{según} [2]). \underline{Entonces el grupo fundamental} $\Pi_1^{(p)}(X)$ \underline{es topológicamente de presentación finita}.

Por SGA 1 IX 4.10, se puede suponer que $k$ es algebraicamente cerrado.


1) \underline{Reducción al caso en que} $X$ \underline{es una superficie regular}.

Si $f : X' \longrightarrow X$ es un morfismo de tipo finito, ponemos $X'' = X'\times_X \sisi{S'}{X'}$ y denotamos por $(X'_a)_{a \in A}$ y $(X''_b)_{b \in B}$ los conjuntos de componentes conexas de $X'$ y $X''$ , respectivamente. Cuando $f$ es un morfismo de descenso efectivo para la categoría de los recubrimientos étales, el grupo fundamental $\Pi_1(X)$ se puede calcular explícitamente
%%%%%%%%%%  scan idx 39  |  folio 28  %%%%%%%%%%
en términos de los grupos fundamentales $\Pi_1(X'_a)$ y $\Pi_1(X''_b)$ ; de ello se deduce (SGA 1 IX 5.2, 5.3) que, si los $\Pi_1(X'_a)$ son topológicamente de generación finita, lo mismo ocurre con $\Pi_1(X)$ ; si los $\Pi_1(X'_a)$ son topológicamente de presentación finita y los $\Pi_1(X''_b)$ topológicamente de generación finita, $\Pi_1(X)$ es topológicamente de presentación finita.

Como $X$ admite una desingularización, se puede encontrar un esquema regular $X'$ y un morfismo propio birracional $f : X' \longrightarrow X$ ; como $f$ es un morfismo de descenso efectivo para la categoría de los recubrimientos étales, basta demostrar el teorema bajo la hipótesis adicional de que $X$ es regular. En efecto, de ello se deduce que $\Pi_1(X)$ es topológicamente de generación finita; sabiendo que $\Pi_1(X')$ es topológicamente de presentación finita y que los $\Pi_1(X''_b)$ son topológicamente de generación finita, se deduce que $\Pi_1(X)$ es topológicamente de presentación finita.

Sea ahora $(U_i)_{i \in I}$ un recubrimiento finito de $X$ por abiertos afines y sea $f : \coprod U_i \longrightarrow X$ el morfismo canónico. Como $f$ es un morfismo de descenso efectivo para la categoría de los recubrimientos étales, se ve que basta demostrar el teorema para los $U_i$ ; se reduce así al caso en que $X$ es afín y liso.

Supongamos en adelante que $X$ es afín, liso sobre $k$ , y que $dim X > 2$ . Se puede aplicar entonces el teorema de Lefschetz para los esquemas cuasiproyectivos ([3] V 7.4); si $Y$ es una sección hiperplana «suficientemente general» de $X$ , se tiene un isomorfismo

\begin{equation*}
\Pi_1(Y) \simeq \Pi_1(X) .
\end{equation*}
Basta, por tanto, demostrar el teorema para $Y$ , lo que reduce al caso en que $dim X = 2$ .


2) \underline{Caso en que} $X$ \underline{es una superficie lisa sobre} $k$

Por el teorema de desingularización fuerte de las superficies, se puede encontrar una superficie $S$ lisa, íntegra y proyectiva y una inmersión abierta $i : X \longrightarrow S$ tal que $Y = S - X$ sea un divisor de cruces normales. Apliquemos el corolario 2.2, cuyas notaciones conservamos.

%%%%%%%%%%  scan idx 40  |  folio 29  %%%%%%%%%%

Mostremos que basta demostrar que el grupo fundamental de $X' = X\times_S S'$ es topológicamente de presentación finita. Como $X$ y $X'$ son regulares y como se tiene

\begin{equation*}
\mathrm{codim}\sisi{(X- \bigcup_{1\leqslant i\leqslant n} \{x_i\},X)}{\left(\bigcup_{1\leqslant i\leqslant n} \{x_i\},X\right)} \geqslant 2 \qquad ,
\end{equation*}
todo recubrimiento étale de $X'$ induce sobre $X- \bigcup_{1\leqslant i\leqslant n}\{x_i\}$ un recubrimiento étale que se prolonga de manera única a un recubrimiento étale de $X$ (SGA 2 1.11); esto muestra que se tiene un isomorfismo

\begin{equation*}
\Pi_1(X') \simeq \Pi_1(X) \qquad .
\end{equation*}
Consideremos el diagrama cartesiano

\begin{tikzcd}
X' \arrow[d, "h"'] & X'_U \arrow[l] \arrow[d, "h_U"] \\
P^1 & U \arrow[l]
\end{tikzcd}
\qquad ,
% STRUCTURE: Commutative square, 4 nodes exactly. Nodes: (1) top-left $X'$ ; (2) top-right $X'_U$ ; (3)
%   bottom-left $P^1$ ; (4) bottom-right $U$. Arrows, 4 exactly, all plain solid arrows with single head
%   (no hook/mono, no iso tilde, no twoheadrightarrow, no dashed): (a) top horizontal, from $X'_U$
%   (top-right) to $X'$ (top-left), pointing LEFT, unlabelled; (b) left vertical, from $X'$ (top-left)
%   down to $P^1$ (bottom-left), pointing DOWN, labelled $h$ with the label printed to the LEFT of the
%   arrow shaft; (c) right vertical, from $X'_U$ (top-right) down to $U$ (bottom-right), pointing DOWN,
%   labelled $h_U$ with the label printed to the RIGHT of the arrow shaft; (d) bottom horizontal, from
%   $U$ (bottom-right) to $P^1$ (bottom-left), pointing LEFT, unlabelled. No cartesian/pullback corner
%   symbol is drawn on the page (the word "cartésien" appears only in the preceding line of prose). A
%   comma is printed on the baseline of the bottom row, well to the right of the node $U$, outside the
%   square.
donde $h = f|X'$ . Denotemos por $L$ el conjunto de los números primos distintos de $p$ y por $\bar{\eta}$ un punto geométrico sobre el punto genérico $\eta$ de $P^1$ . Comprobemos que se satisfacen las hipótesis de SGA 1 XIII 4.8. El morfismo $h$ tiene una sección $\sigma$ porque $f$ la tiene. Demostremos la condición a) de SGA 1 XIII 4.7; como las fibras de $h$ son geométricamente íntegras, el morfismo $h$ es 0-acíclico y localmente 0-acíclico (SGA 4 XV 4.1, 1.16); queda por demostrar que, para todo haz de conjuntos localmente constante $F$ sobre $X'$ , $(F,h)$ es cohomológicamente propio en dimensión $\leqslant 0$ . Si $\xi$ es un punto cerrado de $P^1$ y $Z$ designa la clausura integral de Spec $\mathcal{O}_{P^1,\xi}$ en $\bar{\eta}$ , el esquema $X'_Z = X'\times_{P^1} Z$ satisface la condición $(S_2)$ en los puntos cerrados de $\sisi{X'}{X'_Z}$ y es regular en cualquier otro punto, luego es normal; por consiguiente, un recubrimiento étale de $X'_Z$ cuya restricción a $X'_{\bar{\eta}}$ sea conexa es necesariamente conexo. Esto demuestra nuestra afirmación y, por tanto, a). Como $X'_U$ es el complementario en $S'_U$ de un divisor de cruces normales relativo a $U$ , $h_U$ es localmente 1-asférico para $L$ (SGA 4 XV 2.1) y, para todo haz de $L$-grupos constante y finito $F$ sobre $X'_U$ , $(F,h_U)$ es cohomológicamente propio en dimensión $\leqslant 1$ (SGA 1 XIII 2.4), de donde se sigue la condición b). Como $P^1$ es normal y $T = P^1 - U$ tiene
%%%%%%%%%%  scan idx 41  |  folio 30  %%%%%%%%%%
codimensión 1, se tiene

\begin{equation*}
\text{prof ét}_T(\sisi{S}{P^1}) \geqslant 2 \qquad ,
\end{equation*}
de donde se sigue la condición c). La condición

\begin{equation*}
\text{prof hop}_x(X') \geqslant 3
\end{equation*}
se cumple en todo punto cerrado de $X'$ por el teorema de pureza, de donde se sigue la condición e). Finalmente, para todo punto $t \in P^1$ , $h$ es liso en $\sigma(t)$ , luego es localmente 1-asférico para $L$ en $\sigma(t)$ , de donde se sigue la condición d).

Por SGA 1 XIII 4.7, si $a$ es un punto geométrico de $X'_{\bar{\eta}}$ , se tiene una sucesión exacta

\begin{equation*}
1 \longrightarrow \Pi_1^{(p)}(X'_{\bar{\eta}},a) \longrightarrow \Pi'_1(X'_U,a) \longrightarrow \Pi_1(U,a) \longrightarrow 1 \qquad ,
\end{equation*}
y la elección de la sección $\sigma$ define una escisión de esta sucesión exacta. Denotemos por $K$ la imagen de $\Pi_1(U,a)$ en $\mathrm{Aut}(\Pi_1^{(p)}(X'_{\bar{\eta}},a))$ ; el grupo de coinvariantes bajo $K$ , $\Pi_1^{(p)}(X'_{\bar{\eta}},a)_K$ , es el cociente de $\Pi_1^{(p)}(X'_{\bar{\eta}},a)$ por el subgrupo normal cerrado generado por los elementos del tipo $\sisi{x^{-1}qx}{x^{-1}q(x)}$ , donde $x \in \Pi_1^{(p)}(X'_{\bar{\eta}},a)$  y $q \in K$ . Por loc. cit. (4.7.4), se tiene un isomorfismo

\begin{equation*}
\Pi_1^{(p)}(X',a) \simeq \sisi{\Pi_1^{(p)}(X'_{\bar{\eta}},a)}{\Pi_1^{(p)}(X'_{\bar{\eta}},a)_K} \qquad .
\end{equation*}
Sea $P^1 - U = \coprod_{1\leqslant i\leqslant n} \{t_i\}$ y, para cada $i$ , sea $\bar{T}_i$ la localización estricta de $P^1$ en un punto geométrico $\bar{t}_i$ sobre $t_i$ y sea $s_i$ el punto genérico de $\bar{T}_i$ . Un recubrimiento étale de $U$ cuya restricción a cada $s_i$ es trivial se prolonga a un recubrimiento étale de $P^1$ y, por tanto, es trivial. Para cada $i$ , sea $\bar{k}(s_i)$ la clausura algebraica de $k(s_i)$ ; la afirmación anterior equivale a decir que $\Pi_1(U,a)$ es igual al subgrupo normal cerrado generado por las imágenes de los grupos de Galois $\mathrm{Gal}(\bar{k}(s_i),k(s_i))$ en $\Pi_1(U,a)$ . Por I 5.3, la imagen $K_i$ de $\mathrm{Gal}(\bar{k}(s_i),k(s_i))$ en $\mathrm{Aut}(\Pi_1^{(p)}(X'_{\bar{\eta}},a))$ es finita. Como $K$ es el menor subgrupo normal cerrado que contiene todos los $K_i$ y como $\Pi_1^{(p)}(X'_{\bar{\eta}},a)$ es topológicamente de presentación finita (SGA 1 XIII 3.2), esto demuestra que $\Pi_1^{(p)}(X'_{\bar{\eta}},a)_K$ y, por tanto, también $\Pi_1^{(p)}(X)$, que le es isomorfo, son topológicamente de presentación finita.


%%%%%%%%%%  scan idx 42  |  folio 31  %%%%%%%%%%

\underline{Observación} 2.3.2.  La resolución de singularidades no es necesaria para demostrar el teorema 2.3.1 en el caso en que  $X$  es el abierto complementario de un divisor de cruces normales en un esquema proyectivo liso sobre  $k$ .

\begin{center}\rule{1.5cm}{0.4pt}\end{center}

\underline{BIBLIOGRAFÍA}

[1]\quad S. Abhyankar, \sisi{Local uniformization of algebraic surfaces over ground field of characteristic  $p = 0$ , Amer. J. of Math., t.63, 1956.}{Local uniformization on algebraic surfaces over ground fields of characteristic $p \neq 0$, Ann. of Math., t. 63, no. 3, 1956, p. 491--526.}

[2]\quad H. Hironaka, Resolution of singularities of an algebraic variety over a field of characteristic zero, Ann. of Math., t. 79, 1964.

[3]\quad M. Raynaud, Théorème de Lefschetz en cohomologie des faisceaux cohérents et en cohomologie étale, tesis, de próxima aparición.
